% Generated draft: requires mathematical and source review.
\documentclass[11pt]{article}
\usepackage[margin=25mm]{geometry}
\usepackage{fontspec}
\setmainfont{Latin Modern Roman}
\usepackage{amsmath,amssymb,mathtools,unicode-math}
\setmathfont{Latin Modern Math}
\usepackage{enumitem,needspace,xurl,hyperref}
\setlength{\parindent}{0pt}
\setlength{\parskip}{5pt}
\setlength{\emergencystretch}{3em}
\newcommand{\N}{\mathbb{N}}
\newcommand{\Z}{\mathbb{Z}}
\newcommand{\Q}{\mathbb{Q}}
\newcommand{\R}{\mathbb{R}}
\newcommand{\C}{\mathbb{C}}
\newcommand{\F}{\mathbb{F}}
\newcommand{\A}{\mathbb{A}}
\newcommand{\PP}{\mathbb{P}}
\newcommand{\E}{\mathbb{E}}
\newcommand{\eps}{\varepsilon}
\newcommand{\dd}{\,\mathrm{d}}
\begin{document}
\section*{UnsolvedMath ID: 20002381}
\textbf{Source problem code:} catalogue entry
\subsection{Definitions and mathematical statement}
This participant contribution is a learning goal about Woodin's HOD analysis, not a newly posed open conjecture. Its requested theorem is to understand, under
\[
 L(\R)\models\operatorname{AD},
\]
the analysis of $H=\operatorname{HOD}^{L(\R)}$ and the proof of
\[
 H\models V=\operatorname{HOD},
 \qquad\text{equivalently }\operatorname{HOD}^{H}=H.
\]
Here $\operatorname{HOD}^M$ consists of the hereditarily ordinal-definable sets as computed inside $M$; the second HOD operation is internal to $H$, not a reuse of definitions available in $L(\R)$. The determinacy assumption means that the length-$\omega$ games on natural numbers with payoff sets in $L(\R)$ are determined there. Superscripts lost in the extraction are restored as specified by the supplied source verification.

The supplied research aid provides a general diagnostic criterion rather than the deep HOD analysis. For a transitive model $M\models\operatorname{ZF}$, it reports the equivalence of:
\begin{enumerate}
\item $M\models V=\operatorname{HOD}$;
\item in $M$, there is a set-like global well-order of its universe, definable in the pure membership language from one fixed finite tuple of ordinals of $M$;
\item in $M$, there is a class function $g:\operatorname{Ord}^M\to M$, definable in that same sense, which is total and onto.
\end{enumerate}
``Set-like'' means that the predecessors of each element form a set in $M$; definability refers to a single fixed formula and fixed ordinal parameters, not a different formula or arbitrary set parameter for each element. The surjection statement is a class assertion: a formula $G(\alpha,x,\bar\gamma)$ defines exactly one $x$ for every ordinal $\alpha$ of $M$, and every $x\in M$ is so obtained. It does not assert that the universe or the class function is a set of $M$.

In particular, a well-order or surjection definable only in an expanded structure $(M,\in,P)$ is insufficient by itself. If the extra predicate and the resulting presentation can be eliminated using definitions in $(M,\in)$ from fixed ordinal parameters, the criterion applies. The aid illustrates the distinction with $M=L[c]$ for a Cohen real $c$ over $L$: its canonical $c$-relative well-order is definable using $c$, whereas $c\notin\operatorname{HOD}^{L[c]}$, so $L[c]\not\models V=\operatorname{HOD}$.

Applied to $H=\operatorname{HOD}^{L(\R)}$, the aid's proposed argument is deliberately conditional on the missing deep input that the Woodin--Steel strategy presentation yields such an internally pure-language ordinal-definable surjection. It does not establish that input, and recovering an external or strategy-relative presentation is not by itself an internal proof. The source supplies no precise strategy, recovery theorem, or full bibliographic reference with which to replace that missing hypothesis.
\subsection{Short English statement}
Explain why, assuming determinacy in $L(\R)$, its HOD satisfies $V=\operatorname{HOD}$ internally. The contribution is a learning request; the aid gives a general internal coding criterion and warns that a strategy-relative or externally definable well-order is insufficient.
\subsection{Sources}
\begin{itemize}
\item Peter Koellner, A.9 in the appendix ``Participant Contributions,'' American Institute of Mathematics workshop \emph{Recent advances in core model theory}, \url{https://aimath.org/WWN/coremodel/coremodel.pdf}. The contribution attributes the HOD analysis to Woodin.
\item Local UnsolvedMath registry, permanent ID 20002381, source code \texttt{catalogue entry}, canonical location \texttt{aim-logic-notes.json}, \texttt{notes[156]}, record dated 14 August 2026. The supplied verification identifies missing superscripts; the original literature field is empty.
\item Supplied research attempt 1: context only, problem status at run not a problem, rigor labelled proved, novelty confidence low. Its internal-definability criterion, predicate-elimination observation and Cohen example are machine-generated research aid, not peer reviewed, requiring independent verification before citation. The Woodin--Steel presentation is mentioned without a complete bibliographic reference.
\end{itemize}
\end{document}
